\section*{ВВЕДЕНИЕ}
Технология беспроводной передачи данных Bluetooth, представленная компанией Ericsson в 1994 году, за прошедшие десятилетия стала одним из базовых стандартов организации короткодистанционных соединений. Работая в нелицензируемом диапазоне ISM 2,4 ГГц, она обеспечивает энергоэффективный обмен данными и находит широкое применение в персональных сетях (PAN), системах «умного дома», носимой электронике, а также в устройствах Интернета вещей (IoT). Появление спецификации Bluetooth Low Energy (BLE) существенно расширило возможности технологии в приложениях с требованиями к длительной автономной работе, сделав её одной из ключевых в современных беспроводных решениях.

Несмотря на массовость применения, в учебном и научно-исследовательском процессе технических вузов зачастую ощущается нехватка доступных и гибких аппаратно-программных комплексов. Существующие коммерческие решения либо обладают высокой стоимостью, либо имеют закрытую архитектуру, что ограничивает возможности их модификации и проведения собственных измерений. Большинство готовых лабораторных наборов не позволяют оперативно изменять конфигурацию сети и детально анализировать параметры радиоканала в реальных условиях.

В данной работе объектом исследования является технология беспроводной передачи данных Bluetooth и её реализация в современных сетях передачи данных.

Задачей является разработка модульного учебно-лабораторного стенда, обеспечивающего проведение экспериментов по оценке параметров радиоканала, анализу сетевых топологий и совместимости протокольных стеков, а также разработка лабораторных работ и методических материалов для его использования в учебном процессе.

Предметом исследования является процесс разработки, конфигурирования и применения учебно-лабораторного стенда, включающего автономные узлы на базе микроконтроллеров ESP32 и узлы на базе персональных компьютеров с USB Bluetooth-адаптерами или ESP32, взаимодействующие посредством программного стека BlueZ в среде ОС Linux.

Целью работы является разработка и экспериментальное тестирование учебного лабораторного стенда для исследования технологии Bluetooth в современных беспроводных сетях передачи данных.

В первом разделе работы рассматриваются теоретические основы технологии Bluetooth, анализируются её архитектурные особенности и топологии сетей, а также формулируются технические требования к лабораторному стенду. Во второй главе представлен проект стенда: описывается его структура, обосновывается выбор аппаратной конфигурации узлов, разрабатываются схемы их взаимодействия и способы интеграции с персональными компьютерами. Третья глава посвящена программно-аппаратной реализации стенда, настройке взаимодействия с протокольным стеком BlueZ, а также методикам перехвата и анализа Bluetooth-трафика с использованием программных анализаторов. В четвёртой главе разработан комплект лабораторных работ, включающий методические указания, порядок выполнения экспериментов по исследованию параметров радиоканала, топологий сетей и совместимости протокольных стеков, а также приведены рекомендации по оценке полученных результатов.
